\section{Process tensor theory and construction}
\label{sec:process-tensors}

\markedsubsection{\toctheory}{From quantum channels to multi-time processes}
\label{subsec:from-channels}

% Purpose: move directly from the Liouville-space maps of Section~2 to operations inserted at several times.

% Chronology:
% \begin{itemize}
%     \item Recall the Liouville-space channel \(\lvert\rho'\rangle\!\rangle=\mathsf E\lvert\rho\rangle\!\rangle\).
%     \item Introduce unitary dilation and the CPTP character of a reduced channel for a fixed, initially uncorrelated environment.
%     \item Distinguish deterministic channels from individual CP measurement outcomes.
%     \item Replace uninterrupted evolution by joint system--environment propagators separated by system operations \(\mathcal A_k\).
% \end{itemize}

% Essential equations:
% \begin{align}
%     \mathcal{E}(\rho_S(t_n))
%     &=
%     \operatorname{Tr}_E
%     \bigl[
%         U(\rho_S\otimes\rho_E)U^\dagger
%     \bigr],
%     \\
%     \rho_S(t_n)
%     &=
%     \operatorname{Tr}_E
%     \bigl[
%         U_{n:n-1}(\mathcal A_{n-1}\otimes I_E)
%         \cdots
%         U_{1:0}(\mathcal A_0\otimes I_E)
%         \rho_{SE}(t_0)
%     \bigr],
%     \\
%     \rho_S(t_n)
%     &=
%     \mathcal T_{n:0}[\mathcal A_{n-1},\ldots,\mathcal A_0](\rho_S(t_0)).
% \end{align}

% Figure~1: ``Opening a trajectory''.
% \begin{itemize}
%     \item Top: ordinary joint system--environment evolution.
%     \item Middle: system interventions inserted between propagators.
%     \item Bottom: remove the interventions and leave their legs open, producing \(\mathcal T_{n:0}\).
%     \item Time runs right to left.
% \end{itemize}
Let us zoom out from the tensor-network details of Sec.~\ref{sec:liouville-foundations} for a moment. There, density operators were vectorised, generators were written as many-body MPOs, and their action became an ordinary tensor contraction. Beneath this machinery lies a simple physical requirement: a state evolving from one time to another must remain a valid quantum state. In particular, it must remain Hermitian and positive so that measurement probabilities are real and non-negative, and it must retain unit trace so that those probabilities add up to one. A quantum channel is the linear map that enforces these requirements. Writing
\begin{equation}
    \rho'
    =
    \Phi[\rho],
    \qquad
    \Phi:
    \mathcal B(\mathcal H_{
        \mathrm{in}})
    \longrightarrow
    \mathcal B(\mathcal H_{
        \mathrm{out}}),
\end{equation}
the map \(\Phi\) is required to be completely positive and trace preserving (CPTP),
\begin{align}
    (\Phi\otimes\mathcal I_A)[X] &\geq 0
    &&\text{for every }X\geq0,
    \label{eq:channel-complete-positivity}\\
    \operatorname{Tr}\!\left(\Phi[\rho]\right)
    &=
    \operatorname{Tr}(\rho).
    \label{eq:channel-trace-preservation}
\end{align}
The word \emph{complete} is important. Positivity must be preserved not only for an isolated system, but also when the map acts on one part of a larger state entangled with an otherwise untouched ancilla \(A\). Thus, behind the formal acronym ``CPTP'' are the familiar probability rules of quantum mechanics, together with the demand that they remain consistent when the system is embedded in a larger experiment~\cite{Taranto2025}.

We have already encountered one such example, although we did not call it a channel at the time. In Sec.~\ref{sec:liouville-foundations}, Markovian open-system dynamics was generated by the Gorini--Kossakowski--Lindblad--Sudarshan (GKLS) equation
\begin{equation}
    \frac{\mathrm d\rho}{\mathrm dt}
    =
    \mathcal L_t[\rho]
    =
    -\mathrm i[H(t),\rho]
    +
    \sum_{\mu}\gamma_{\mu}(t)
    \left(
        L_{\mu}\rho L_{\mu}^{\dagger}
        -
        \frac{1}{2}
        \left\{
            L_{\mu}^{\dagger}L_{\mu},\rho
        \right\}
    \right),
    \qquad
    \gamma_{\mu}(t)\geq0.
    \label{eq:gkls-generator-channel}
\end{equation}
Strictly speaking, the Liouvillian \(\mathcal L_t\) is the generator rather than the channel itself. The corresponding propagator
\begin{equation}
    \Phi_{t:s}
    =
    \overleftarrow{\mathcal T}
    \exp\!\left[
        \int_s^t \mathrm d\tau\,\mathcal L_{\tau}
    \right],
    \qquad
    \rho(t)=\Phi_{t:s}[\rho(s)],
    \label{eq:gkls-propagator-channel}
\end{equation}
is CPTP; for a time-independent generator this reduces to \(\Phi_{t:s}=\exp[(t-s)\mathcal L]\). The Liouvillian MPO introduced previously is therefore the tensor-network representation of the generator whose propagation carries one density operator to another without leaving the space of physical states.

A second, and more general, route to a quantum channel is to regard the apparent non-unitarity of the system as part of a larger unitary evolution. It is straightforward to see that a unitary transformation on a many-body denstiy operator is a CPTP map: they preserve the state norm by definition, \(U\rho U^\dagger\) is Hermitian if \(\rho\) is Hermitian, and it keeps the same eigenvalues under the transformation, ensuring positivity. If the system begins uncorrelated with an environment in a fixed state \(\rho_E\), then
\begin{equation}
    \Phi_{t:s}[\rho_S]
    =
    \operatorname{Tr}_E
    \left[
        U_{t:s}
        (\rho_S\otimes\rho_E)
        U_{t:s}^{\dagger}
    \right]
    =
    \sum_{\mu}
    K_{\mu}\rho_S K_{\mu}^{\dagger},
    \label{eq:channel-dilation-kraus}
\end{equation}
where the Kraus operators obey
\begin{equation}
    \sum_{\mu}K_{\mu}^{\dagger}K_{\mu}
    =
    I_S.
    \label{eq:kraus-completeness}
\end{equation}
In finite dimensions, a linear map is CPTP if and only if it admits such a Kraus representation; equivalently, every quantum channel can be realised by adjoining an initially uncorrelated environment, evolving the joint state unitarily, and discarding the environment afterwards~\cite{Taranto2025}. Unitary evolution and GKLS propagation are therefore not separate ideas, but familiar members of the same family of physical state-to-state maps.

All these channels answer essentially the same question: given the state of the system now, what will its state be later? For a memoryless evolution, this state-to-state picture may be applied successively,
\begin{equation}
    \rho_S(t_{k+1})
    =
    \Phi_{k+1:k}[\rho_S(t_k)].
    \label{eq:stepwise-channel-picture}
\end{equation}
But must the map \(\Phi_{k+1:k}\) be determined by the present reduced state alone? Suppose two experimental histories produce exactly the same \(\rho_S(t_k)\), while leaving behind different correlations between the system and its environment. An operation performed at an earlier time may have disturbed the environment, conditioned it on a measurement outcome, or stored information in degrees of freedom that are no longer visible in \(\rho_S(t_k)\). The same current system state can then have different futures. In that situation, the question ``what state comes next?'' has omitted part of the experimental input: ``how did we reach the present state?''.

The natural remedy is to promote the history of interventions to an input of the dynamical description. Let \(\mathcal A_k\) denote the operation applied to the system at time \(t_k\), while the joint system--environment dynamics between consecutive times is described by \(\mathcal U_{k+1:k}\). Starting from a possibly correlated state \(\rho_{SE}(t_0)\), the final reduced state is
\begin{align}
    \rho_S(t_n)
    ={}&
    \operatorname{Tr}_E
    \Bigl[
        \mathcal U_{n:n-1}
        (\mathcal A_{n-1}\otimes\mathcal I_E)
        \cdots
        \mathcal U_{1:0}
        (\mathcal A_0\otimes\mathcal I_E)
        [\rho_{SE}(t_0)]
    \Bigr]
    \nonumber\\
    =:{}&
    \mathcal T_{n:0}
    \bigl[
        \mathcal A_{n-1},\ldots,\mathcal A_0
    \bigr].
    \label{eq:process-tensor-operational-definition}
\end{align}
The process tensor \(\mathcal T_{n:0}\) is the higher-order map that remains after the initial joint state, the inaccessible environment, and all intermediate system--environment propagators have been absorbed into a single object~\cite{Pollock2018Framework,Milz2021,Taranto2025}. A channel accepts a state and returns a state; a one-slot superchannel accepts an operation and returns a state or another operation. A process tensor (also referred to as quantum combs, commonly in the literature) is their multi-time extension: it accepts an ordered sequence of operations and returns the final state. If the final output is also measured or traced, the same object instead returns the associated probability or expectation value.

\begin{figure}[t]
    \centering
    \fbox{%
        \parbox{0.91\linewidth}{
            \centering
            \vspace{0.6em}
            \textbf{Figure placeholder: from a channel to a multi-time process}\par
            \vspace{0.5em}
            Panel (a): a single CPTP map \(\Phi_{t:s}\) taking
            \(\rho_S(s)\) to \(\rho_S(t)\).\par
            \vspace{0.35em}
            Panel (b): joint system--environment propagation interrupted by
            \(\mathcal A_0,\ldots,\mathcal A_{n-1}\), with the environment
            carrying information between the intervention times.\par
            \vspace{0.35em}
            Panel (c): the same microscopic network enclosed inside
            \(\mathcal T_{n:0}\), leaving one open slot for each
            \(\mathcal A_k\) and a final system output. Time should run from
            right to left, consistently with the remaining tensor diagrams.
            \vspace{0.6em}
        }
    }
    \caption{
        A quantum channel describes one transformation from an input state to
        an output state. A process tensor retains the underlying
        system--environment evolution while leaving the intermediate system
        operations open, thereby describing an entire multi-time experiment.
    }
    \label{fig:channel-to-multitime-process}
\end{figure}

Despite acting on operations rather than directly on states, a physical process tensor inherits three familiar structural properties from quantum mechanics~\cite{Pollock2018Framework,Milz2021,Taranto2025}. First, it is \emph{linear in every intervention}. If an operation at time \(t_k\) is chosen as a probabilistic mixture \(p\mathcal A_k+(1-p)\mathcal B_k\), then the resulting output is the same mixture of the outputs obtained from \(\mathcal A_k\) and \(\mathcal B_k\). For example,
\begin{align}
    &\mathcal T_{n:0}
    \bigl[
        \ldots,
        p\mathcal A_k+(1-p)\mathcal B_k,
        \ldots
    \bigr]
    \nonumber\\
    &\qquad=
    p\,
    \mathcal T_{n:0}
    \bigl[\ldots,\mathcal A_k,\ldots\bigr]
    +(1-p)\,
    \mathcal T_{n:0}
    \bigl[\ldots,\mathcal B_k,\ldots\bigr].
    \label{eq:process-tensor-linearity}
\end{align}

Second, the process tensor is \emph{completely positive}. Inserting completely positive operations at its open times must produce a positive final state, even when those operations act jointly on the system and an arbitrary untouched ancilla. This is the multi-time counterpart of Eq.~\eqref{eq:channel-complete-positivity}. It guarantees that the process remains physically meaningful when it is embedded inside a larger circuit, rather than only for the isolated sequence used to define it.

Finally, the process tensor is \emph{causal}. A choice made at a later time cannot alter statistics already recorded at an earlier time. Operationally, once the future slots are filled with deterministic trace-preserving operations and their outputs are discarded, the probabilities assigned to the earlier interventions cannot depend on which future operations were chosen. This no-signalling-from-the-future condition also supplies the appropriate multi-time notion of normalisation: every deterministic sequence must preserve total probability. In the Choi representation it becomes a hierarchy of partial-trace constraints, which we introduce in the following subsection. Together, linearity, complete positivity, and causality distinguish a physical time-ordered process from an arbitrary multilinear tensor with matching dimensions.

A point worth emphasising is that an individual trajectory selected from a process tensor need not be deterministic. An operation inserted at an intermediate time may represent a probabilistic event, which would then, reduce the trace of the outcome quantum state. Unit trace is determined by the choice of instruments, and so, trace-preserving trajectory is not a requirement for a process tensor. We will see more on this in section 4 (Quantum instruments) \textcolor{red}{TODO: add link}.

%%%%%%%%%%%%%%%%%%%%%%%%%%%%%%%%%%%%%%%%%%%%%%%%
\markedsubsection{\toctheory}{Choi representation, instruments, and process contraction}
\label{subsec:choi-instruments}

% Purpose: provide exactly enough formalism to explain what is stored and how it is evaluated.

% Chronology:
% \begin{itemize}
%     \item Recall the channel--state correspondence.
%     \item Apply it at every temporal slot to obtain the process Choi operator \(\Upsilon_{n:0}\).
%     \item Identify the input and output space belonging to each intervention time.
%     \item Introduce quantum instruments as collections of CP operations,
%     \[
%         \mathcal J_k
%         =
%         \bigl\{\mathcal A_k^{(x_k)}\bigr\}_{x_k},
%         \qquad
%         \sum_{x_k}\mathcal A_k^{(x_k)}
%         \ \text{is CPTP}.
%     \]
%     \item State the generalized Born rule.
%     \item Explain the three possible results of a contraction:
%     \begin{itemize}
%         \item an open final leg gives a state;
%         \item a completely closed network gives a probability or expectation value;
%         \item selecting measurement outcomes gives an unnormalized conditional state.
%     \end{itemize}
% \end{itemize}

% Essential equations:
% \begin{align}
%     J_{\mathcal A}
%     &=
%     (\mathcal A\otimes I)\bigl[\lvert\Omega\rangle\langle\Omega\rvert\bigr],
%     \\
%     \rho_S(t_n)
%     &=
%     \Upsilon_{n:0}
%     \ast J_{\mathcal A_{n-1}}
%     \ast\cdots\ast J_{\mathcal A_0},
%     \\
%     p(x)
%     &=
%     \operatorname{Tr}
%     \Bigl[
%         \Upsilon_{n:0}
%         \Bigl(
%             \bigotimes_{k=0}^{n-1}
%             J_{\mathcal A_k^{(x_k)}}
%         \Bigr)^{T_{\mathrm{slots}}}
%     \Bigr].
% \end{align}
% The transpose/link-product convention must be derived from the column-major convention fixed in Section~2.

% Figure~2: ``The Born rule across time''.
% \begin{itemize}
%     \item State + effect \(\to\) probability.
%     \item Channel Choi tensor + state \(\to\) output state.
%     \item Process Choi tensor + instrument sequence \(\to\) conditional state or probability.
% \end{itemize}

% Package connection: the diagram and code should perform the identical contraction.
% The constructor names below can be replaced by the final public API.
% \begin{verbatim}
% seq = InstrumentSeq(pt)
% add!(seq, 0, prepare(rho0))
% add!(seq, k, measure(M))
% add!(seq, nsteps, trace_out())
% result = evaluate_process(pt, seq)
% \end{verbatim}

%%%%%%%%%%%%%%%%%%%%%%%%%%%%%%%%%%%%%%%%%%%%%%%%
\markedsubsection{\tocboth}{Anatomy of a \texttt{ProcessTensor}}
\label{subsec:pt-anatomy}

% Purpose: merge the remaining theory with the concrete object stored by the package.

% Topics:
% \begin{itemize}
%     \item Represent \(\Upsilon_{n:0}\) as an MPO ordered along time.
%     \item Explain one PT-MPO core before showing the complete chain.
%     \item Identify:
%     \begin{itemize}
%         \item temporal input legs;
%         \item temporal output legs;
%         \item memory/link bonds;
%         \item initial and final boundary tensors;
%         \item coupling times;
%         \item any final open system leg.
%     \end{itemize}
%     \item Explain that physical validity requires positivity, normalization, and causal ordering.
%     \item Give the recursive causality condition without a separate derivation:
%     \[
%         \operatorname{Tr}_{o_k}\Upsilon_{k:0}
%         =
%         I_{i_k}\otimes\Upsilon_{k-1:0}.
%     \]
%     \item Mention the Markovian factorization only as a limiting anatomy:
%     \[
%         \Upsilon_{n:0}^{\mathrm M}
%         =
%         \rho_0
%         \otimes J_{\Phi_{1:0}}
%         \otimes\cdots\otimes
%         J_{\Phi_{n:n-1}}.
%     \]
%     \item Interpret nontrivial MPO bonds as carriers of temporal correlations, while avoiding detailed memory measures here.
% \end{itemize}

% Figure~3: ``Anatomy of the stored process tensor''.
% A single enlarged PT-MPO with callouts for:
% \begin{itemize}
%     \item \texttt{input\_sites(pt)};
%     \item \texttt{output\_sites(pt)};
%     \item \texttt{coupling\_times(pt)};
%     \item \texttt{bond-\%02d};
%     \item open and contracted boundary legs;
%     \item right-to-left time direction.
% \end{itemize}

% Inspection code:
% \begin{verbatim}
% input_sites(pt)
% output_sites(pt)
% coupling_times(pt)
% open_leg_info(pt)
% isfullycontracted(pt)
% \end{verbatim}

% This subsection should make the transition: we now know what the package must return. The remaining question is what physical information must be supplied to construct it. That is the natural 50--50 turning point.

The Choi representation introduced above is complete, but its size grows exponentially with the number of intervention times. This is the temporal counterpart of the many-body problem encountered in Sec.~\ref{sec:liouville-foundations}: writing down every possible history as one dense tensor quickly becomes impractical. The same matrix-product architecture that compressed spatial many-body states can instead be turned along the time direction. Splitting the temporal indices into consecutive groups gives
\begin{equation}
    \bigl[\mathcal T_{n:0}\bigr]_%
    {\boldsymbol{o},\boldsymbol{i}}
    =
    \sum_{\alpha_1,\ldots,\alpha_{n-1}}
    Q^{[n-1],o_{n-1}i_{n-1}}_{\alpha_{n-1}}
    Q^{[n-2],o_{n-2}i_{n-2}}_{\alpha_{n-2}\alpha_{n-1}}
    \cdots
    Q^{[0],o_0i_0}_{\alpha_1}.
    \label{eq:pt-mpo-decomposition}
\end{equation}
Here \(i_k\) and \(o_k\) are the system input and output indices at time step \(k\), while \(\alpha_k\) is an internal memory index joining neighbouring cores. The dimensions \(\chi_k=\dim(\alpha_k)\) quantify how large an effective space must be carried across each temporal cut. A Markovian process factorises across time and needs only trivial memory links, whereas a non-Markovian process generally has \(\chi_k>1\). A small value of \(\chi_k\) does not imply the absence of memory; it means that the relevant memory is compressible~\cite{Pollock2018Framework,Jorgensen2019}.

Equation~\eqref{eq:pt-mpo-decomposition} is stored by \texttt{ProcessTensors.jl} as a Liouville-space MPO whose sites are ordered in time. An interior core has two temporal legs and two memory links. The input leg is primed, with prime level \(1\), and the output leg is unprimed, with prime level \(0\). Both have dimension \(d_S^2\), retain the physical metadata of the system Liouville site, and carry a \texttt{tstep=k} tag. The memory links carry the tags \texttt{PT}, \texttt{Link}, and their temporal position. The first and last bath links are absent from the final MPO because they have already been contracted with the initial bath state \(\lvert\rho_E(0)\rangle\rangle\) and the final trace \(\langle\!\langle I_E\rvert\), respectively.

\begin{figure}[t]
    \centering
    \fbox{%
        \parbox{0.91\linewidth}{
            \centering
            \vspace{0.6em}
            \textbf{Figure placeholder: anatomy of the temporal MPO}\par
            \vspace{0.5em}
            Draw four rounded PT cores with time running from right to left.
            Each core \(Q^{[k]}\) has a lower primed input leg
            \(i_k\), an upper unprimed output leg \(o_k\), and thick
            horizontal memory links \(\alpha_k\). Label the physical legs
            by their ITensor data: dimension \(d_S^2\), prime level, and
            \texttt{tstep=k}. Label the internal legs
            \texttt{PT,Link,tstep=k}. At the two boundaries, show
            \(\lvert\rho_E(0)\rangle\rangle\) and
            \(\langle\!\langle I_E\rvert\) being absorbed into the first
            and last cores.
            \vspace{0.6em}
        }
    }
    \caption{
        Matrix-product representation of a process tensor. Open system legs
        mark the times at which operations may be inserted, while the
        horizontal bonds transmit the compressed influence of the environment
        between different parts of the experiment.
    }
    \label{fig:process-tensor-anatomy}
\end{figure}

The package wraps this MPO together with the physical information used to construct it. Concretely, a \texttt{ProcessTensor} stores the following fields:
\begin{center}
\begin{tabular}{|p{0.20\linewidth}|p{0.70\linewidth}|}
    \hline
    \textbf{Field} & \textbf{Meaning} \\
    \hline
    \texttt{core} & The time-ordered \texttt{MPO\{Liouville\}} containing the tensors \(Q^{[k]}\). \\
    \hline
    \texttt{system} & The system Hamiltonian, optional Lindblad channels, and canonical Liouville site indices. \\
    \hline
    \texttt{environment} & The bath modes, their initial states and Hamiltonians, and the system--bath couplings. It is \texttt{nothing} for a bath-free process. \\
    \hline
    \texttt{dt}, \texttt{nsteps} & The time spacing and number of temporal MPO cores. The represented final time is \(t_{\mathrm f}=n_{\mathrm{steps}}\Delta t\). \\
    \hline
    \texttt{coupling\_site} & The canonical system Liouville index whose copies become the input and output legs of every PT core. \\
    \hline
\end{tabular}
\end{center}
Unknown property access is delegated to \texttt{core}, so familiar tensor-network inspection functions can be applied directly to \texttt{pt}. For example, \texttt{linkdims(pt)} and \texttt{maxlinkdim(pt)} inspect its temporal memory bonds, while \texttt{input\_sites}, \texttt{output\_sites}, and \texttt{coupling\_times} return the exact ITensor indices needed at particular times. Reusing these returned indices is essential: ITensor contractions match index identity, not merely equal dimensions or similar-looking tags.

%%%%%%%%%%%%%%%%%%%%%%%%%%%%%%%%%%%%%%%%%%%%%%%%
\markedsubsection{\tocjulia}{Ingredients of a process-tensor construction in Julia}
\label{subsec:pt-ingredients}

% Purpose: introduce the builder inputs one at a time, physically and computationally.
% Use one running example throughout---a single spin coupled to a small spin environment---so the reader watches one process tensor being assembled.

% \paragraph{1. System space.}
% Introduce Hilbert-space system indices, corresponding Liouville indices, the initial system state, and the system Hamiltonian or Liouvillian.
% Clarify whether the reusable process contains the free system dynamics or represents only the environmental influence; this must match the package's exact convention.
% \begin{verbatim}
% system_sites = siteinds("S=1/2", 1)
% system_liouv = liouv_sites(system_sites)
% rhoS0 = ...
% HS    = ...
% \end{verbatim}

% \paragraph{2. Environment.}
% Introduce bath indices, bath Hamiltonian, initial bath state, product versus correlated bath initialization, and why the bath can be discarded after its influence has been encoded.
% \begin{verbatim}
% bath_sites = siteinds("S=1/2", Nb)
% bath_liouv = liouv_sites(bath_sites)
% rhoE0 = ...
% HE    = ...
% \end{verbatim}

% \paragraph{3. System--environment coupling.}
% Introduce the decomposition
% \[
%     H_{SE}
%     =
%     \sum_{\alpha} S_{\alpha}\otimes B_{\alpha},
% \]
% and explain that the coupling operators determine which bath information is visible through the system.
% \begin{verbatim}
% couplings = [
%     (Salpha, Balpha),
%     (Sbeta,  Bbeta),
% ]
% \end{verbatim}

% \paragraph{4. Temporal discretization.}
% Introduce
% \[
%     t_k=t_0+k\Delta t,\qquad k=0,\ldots,n,
% \]
% together with \texttt{dt}, \texttt{nsteps}, \texttt{coupling\_times}, and the distinction between time-discretization error and tensor-network truncation.

% \paragraph{5. Construction controls.}
% Briefly introduce cutoff or truncation tolerance, maximum bond dimension, contraction/compression schedule, and builder or solver choice.
% Do not explain ACE or TEMPO internally yet. These are controls shared by builders; the algorithms belong in later sections.

% Figure~4: ``Ingredients entering a builder''.
% A compact mapping rather than a large flowchart:

% \begin{center}
% \begin{tabular}{@{}lll@{}}
% \hline
% Physical ingredient & Mathematical object & Package object \\
% \hline
% System & \(\mathcal H_S\), \(H_S\), \(\rho_S\) & system indices/model \\
% Environment & \(\mathcal H_E\), \(H_E\), \(\rho_E\) & bath or modes \\
% Coupling & \(\sum_{\alpha}S_{\alpha}\otimes B_{\alpha}\) & coupling operators \\
% Time grid & \(\Delta t\), \(n\) & \texttt{dt}, \texttt{nsteps} \\
% Approximation & temporal compression & cutoff, max bond \\
% \hline
% \end{tabular}
% \end{center}

We now construct the smallest process that can carry non-trivial memory: one system spin coupled to one bath spin. Keeping the example deliberately small lets the dense builder retain the complete bath Liouville space, so every object printed by the package can be interpreted directly. We begin with the imports
\begin{juliacode}[
  caption={Imports used throughout the dense process-tensor example.},
  label={lst:pt-ingredients-imports},
]
using ITensors
using ProcessTensors
using ITensors.Ops: Exact, Trotter
\end{juliacode}
and introduce the required ingredients one at a time.

The first ingredient is the accessible system. Let
\begin{equation}
    H_S=hS_S^x,
    \qquad h=0.7.
    \label{eq:dense-example-system-hamiltonian}
\end{equation}
It is constructed from an ordinary Hilbert-space spin index:
\begin{juliacode}[
  caption={System Hamiltonian and \texttt{SpinSystem} constructor.},
  label={lst:pt-ingredients-system},
]
system_sites = siteinds("S=1/2", 1)

h = 0.7
H_S = OpSum()
H_S += h, "Sx", 1

system = spin_system(system_sites, H_S)
println(system)
\end{juliacode}
which prints
\begin{juliaoutput}
ProcessTensors.SpinSystem
  sites: 1
  space: Liouville
  site dims: 4
  dissipative: false
\end{juliaoutput}
The first line identifies the model type. Although we supplied a Hilbert index of dimension \(2\), \texttt{spin\_system} converted it internally into a Liouville index of dimension \(2^2=4\); this accounts for the next three lines. The final line reports that no Lindblad jump operators were supplied. The present process-tensor builder requires a single-site system because the free-system propagation is embedded into every temporal core. An initial system state is intentionally absent: it is an experimental preparation and will later be inserted into the first open process-tensor leg as an instrument.

The inaccessible environment requires more information because it is absorbed into the process tensor rather than supplied later. For each bath mode we need its Liouville site, local Hamiltonian, initial state, and coupling to the system. We take
\begin{equation}
    H_E=\frac{\omega_E}{2}S_E^z,
    \qquad
    H_{SE}=gS_E^zS_S^z,
    \qquad
    \rho_E(0)=\lvert\!\downarrow\rangle
              \langle\downarrow\!\rvert,
    \label{eq:dense-example-bath-model}
\end{equation}
with \(\omega_E=1.1\) and \(g=0.35\):
\begin{juliacode}[
  caption={Bath Hamiltonian, initial state, coupling, and \texttt{SpinBath} constructor.},
  label={lst:pt-ingredients-bath},
]
bath_sites = siteinds("S=1/2", 1)
bath_sites_L = liouv_sites(bath_sites)

omega_E = 1.1
H_E = OpSum()
H_E += (omega_E / 2), "Sz", 1

psi_E0 = MPS(bath_sites, ["Dn"])
rho_E0 = to_dm(psi_E0)
rho_E0_L = to_liouville(rho_E0; sites=bath_sites_L)

g = 0.35
H_SE = OpSum()
H_SE += g, "Sz", 1, "Sz", 2

mode = spin_mode(
    bath_sites_L,
    H_E,
    rho_E0_L;
    coupling=H_SE,
)
environment = spin_bath([mode])

println(environment)
\end{juliacode}
The two site labels in \(H_{SE}\) follow a local convention used by every independent bath mode: site \(1\) is the bath mode and site \(2\) is the system coupling site. They do not refer to two sites inside the system. The call above prints
\begin{juliaoutput}
ProcessTensors.SpinBath
  modes: 1
  space: Liouville
  site dims: 4
  bath Liouville dimension: 4

  mode summary:
    [1] SpinMode(dim=4, H_terms=1, coupling_terms=1)
\end{juliaoutput}
This output first identifies a spin bath containing one mode. Its single Liouville site again has dimension \(4\), so the complete bath Liouville space also has dimension \(D_E=4\). The mode summary confirms that both the local bath Hamiltonian and the mode--system coupling contain one \texttt{OpSum} term. Notice that \(\rho_E(0)\) was vectorised on the exact indices stored in \texttt{bath\_sites\_L}. Two Liouville indices with the same dimension and tags are not interchangeable if they have different ITensor identities.

The last physical ingredient is the time grid,
\begin{equation}
    t_k=k\Delta t,
    \qquad
    k=0,\ldots,n_{\mathrm{steps}},
    \label{eq:process-tensor-time-grid}
\end{equation}
which is specified by
\begin{juliacode}[
  caption={Time grid for the dense process-tensor example.},
  label={lst:pt-ingredients-time},
]
dt = 0.05
nsteps = 5
println("time grid: dt=$dt, nsteps=$nsteps, t_final=$(dt * nsteps)")
\end{juliacode}
and produces
\begin{juliaoutput}
time grid: dt=0.05, nsteps=5, t_final=0.25
\end{juliaoutput}
The builder returns five PT-MPO cores, one for each interval of duration \(\Delta t\). Their time tags run from \texttt{tstep=0} to \texttt{tstep=4}; the terminal boundary lies at \(t_5=0.25\).

%%%%%%%%%%%%%%%%%%%%%%%%%%%%%%%%%%%%%%%%%%%%%%%%
\markedsubsection{\tocjulia}{Building and inspecting the process tensor}
\label{subsec:build-inspect}

% Purpose: bring the ingredients together into the first complete construction.

% The presentation should follow exactly this rhythm:
% \begin{enumerate}
%     \item Equation for one joint time step,
%     \[
%         U_{k+1:k}
%         =
%         e^{\Delta t\,\mathcal L_{SE}(t_k)}.
%     \]
%     \item Diagram showing the bath legs joined between steps while the system operation slots remain open.
%     \item Code building the same network.
%     \item Diagram showing the resulting \texttt{ProcessTensor}.
% \end{enumerate}

% Schematic code:
% \begin{juliacode}[
%   caption={Schematic construction of a process tensor from the ingredients above.},
%   label={lst:pt-build-schematic},
% ]
% environment = ...
% system      = ...
% schedule    = default_schedule(dt, nsteps)

% pt = build_process_tensor(
%     system,
%     coupling_site;
%     environment,
%     dt,
%     nsteps,
%     cutoff,
%     maxdim,
% )
% \end{juliacode}

% Then inspect rather than immediately use it:
% \begin{juliacode}[
%   caption={Inspection of the stored process-tensor MPO.},
%   label={lst:pt-inspect-schematic},
% ]
% input_sites(pt)
% output_sites(pt)
% coupling_times(pt)
% open_leg_info(pt)
% maxlinkdim(pt)
% \end{juliacode}

% Figure~5: ``From microscopic model to \texttt{ProcessTensor}''.
% \begin{itemize}
%     \item Left: system, bath, initial state, coupling, and time grid.
%     \item Centre: open system--bath evolution network.
%     \item Right: compressed PT-MPO returned by the builder.
%     \item Use identical colors and index tags in all three panels.
% \end{itemize}

% This is where the reader should finally be able to say what every argument to \texttt{build\_process\_tensor} means.


For this first construction we use the \texttt{Dense()} backend. It forms the one-step joint Liouville propagator on the full system--bath space and reshapes the bath input and output indices into adjacent temporal memory links. Suppressing the free-system splitting for a moment, a single core has the schematic form
\begin{equation}
    Q^{[k],o_ki_k}_{\alpha_{k+1}\alpha_k}
    =
    \langle\!\langle\alpha_{k+1}\rvert
    \exp\!\left(\Delta t\,\mathcal L_{SE}\right)
    \lvert\alpha_k\rangle\!\rangle,
    \label{eq:dense-pt-core}
\end{equation}
where \(\alpha_k\) runs over a basis of the bath Liouville space. The initial boundary is contracted with \(\lvert\rho_E(0)\rangle\rangle\), the final boundary with \(\langle\!\langle I_E\rvert\), and the intermediate bath legs are identified between consecutive cores. For a single bath of Hilbert dimension \(d_E\), the uncompressed memory link can therefore have dimension \(d_E^2\). This direct construction is transparent and numerically exact for the chosen time-step propagators, but its joint matrix dimension grows as \(d_S^2d_E^2\); it is consequently intended for no bath, one mode, or small multimode reference calculations.

The complete call for our example is
\begin{juliacode}[
  caption={Dense construction of the example process tensor.},
  label={lst:pt-build-dense},
]
pt = build_process_tensor(
    system;
    method=Dense(),
    environment=environment,
    dt=dt,
    nsteps=nsteps,
    alg=Exact(),
    sys_alg=Trotter{2}(),
    progress=false,
    verbose=true,
)
\end{juliacode}
The keyword \texttt{method=Dense()} selects the process-tensor construction backend. By contrast, \texttt{alg=Exact()} specifies how the joint bath-plus-coupling Liouville propagator is exponentiated. The system's free one-step map is then embedded into every bath core. The symmetric choice \texttt{sys\_alg=Trotter\{2\}()} places half of this free-system map on either side of the bath core,
\begin{equation}
    M_S(\Delta t/2)\,
    Q_{E+SE}(\Delta t)\,
    M_S(\Delta t/2),
    \label{eq:symmetric-system-bath-sandwich}
\end{equation}
whereas the default \texttt{Trotter\{1\}()} uses the asymmetric ordering \(Q_{E+SE}(\Delta t)M_S(\Delta t)\). Thus \texttt{Dense()} and \texttt{Exact()} answer different questions: the former chooses how the multi-time object is assembled, while the latter chooses how one local propagator is computed.

We have disabled the transient progress animation and enabled persistent logging so that the stages remain visible after the calculation. Omitting terminal-dependent decorations, the output reads
\begin{juliaoutput}
Info: Building process tensor
  method = Dense
  nsteps = 5
  dt = 0.05
  environment = SpinBath(nmodes=1, D_bath=4, coupling=true)

Info: Constructing joint propagator
  joint_dimension = 16
  channel = liouvillian

Info: Assembling process-tensor cores
  total = 5

Info: Closing bath boundaries

Info: Built process tensor
  maxlinkdim = 4
  elapsed_seconds = ...
\end{juliaoutput}
The opening record repeats the selected backend, temporal grid, and bath summary. The joint dimension \(16=4\times4\) is the product of the system and bath Liouville dimensions; this is the matrix space on which the dense exponential is constructed. The next stage places one copy of that propagator into each of the five temporal slabs. ``Closing bath boundaries'' means contracting the first bath link with \(\rho_E(0)\) and the last with the trace effect. Finally, \texttt{maxlinkdim=4} reports the largest surviving temporal memory space. The elapsed time is deliberately machine dependent.

The returned object has its own compact summary:
\begin{juliacode}[
  caption={Printed summary of the constructed process tensor.},
  label={lst:pt-println},
]
println(pt)
\end{juliacode}
\begin{juliaoutput}
5-step ProcessTensor{SpinSystem, SpinBath} | dt=0.05 | t_final=0.25 | maxlinkdim=4

  system:      SpinSystem(nsites=1, dissipative=false)
  environment: SpinBath(nmodes=1, D_bath=4, coupling=true)
  core:        MPO{Liouville}(length=5, linkdims=[4, 4, 4, 4])
\end{juliaoutput}
The header confirms the number of temporal cores, their spacing, the represented final time, and the largest memory bond. The system and environment lines record the physical objects absorbed into the construction. The final line exposes the tensor-network anatomy: five Liouville MPO cores joined by four memory links, each of dimension \(4\). No temporal compression has been performed by \texttt{Dense()}, so these links retain the complete one-spin bath Liouville space.

We can inspect the open system legs without relying on the randomly generated textual IDs of individual ITensor indices:
\begin{juliacode}[
  caption={Inspection of open system legs and temporal memory bonds.},
  label={lst:pt-inspect-legs},
]
in0 = only(input_sites(pt, 0))
out0 = only(output_sites(pt, 0))
tin = tag_value(in0, "tstep=")
tout = tag_value(out0, "tstep=")

println("input : dim=$(dim(in0)), plev=$(plev(in0)), tstep=$tin")
println("output: dim=$(dim(out0)), plev=$(plev(out0)), tstep=$tout")
println("memory bond dimensions: ", linkdims(pt))
println("largest memory bond: ", maxlinkdim(pt))

out_prev, in_curr = coupling_times(pt, 1)
tprev = tag_value(only(out_prev), "tstep=")
tcurr = tag_value(only(in_curr), "tstep=")
println("step 1 connects output tstep=$tprev to input tstep=$tcurr")
\end{juliacode}
The resulting output is
\begin{juliaoutput}
input : dim=4, plev=1, tstep=0
output: dim=4, plev=0, tstep=0
memory bond dimensions: [4, 4, 4, 4]
largest memory bond: 4
step 1 connects output tstep=0 to input tstep=1
\end{juliaoutput}
The equal physical dimensions show that both legs represent the same system Liouville space; their prime levels distinguish the state entering a temporal slab from the state leaving it. The \texttt{tstep} tags locate those legs within the process. Finally, \texttt{coupling\_times(pt, 1)} returns the output at step \(0\) and the input at step \(1\), which are precisely the two legs on which an intermediate two-leg operation will act.

\begin{figure}[t]
    \centering
    \fbox{%
        \parbox{0.91\linewidth}{
            \centering
            \vspace{0.6em}
            \textbf{Figure placeholder: dense construction in equation,
            diagram, and code}\par
            \vspace{0.5em}
            Panel (a): the one-spin system, one-spin bath, their initial
            data, and \(H_{SE}\). Panel (b): five copies of the dense joint
            propagator, with bath legs connected horizontally, the first bath
            leg closed by \(\rho_E(0)\), and the last by
            \(\langle\!\langle I_E\rvert\). Panel (c): the resulting
            \texttt{ProcessTensor} summary and the corresponding
            \texttt{build\_process\_tensor} call. Keep the same index colors
            and tags used in Fig.~\ref{fig:process-tensor-anatomy}.
            \vspace{0.6em}
        }
    }
    \caption{
        Dense construction of a process tensor from its microscopic
        ingredients. The inaccessible bath is closed at the temporal
        boundaries but retained between time steps as the MPO memory bond;
        the system legs remain open for later interventions.
    }
    \label{fig:dense-process-tensor-construction}
\end{figure}

At this stage the process has been constructed and inspected, but no experimental question has yet been asked of it. The open legs in Fig.~\ref{fig:dense-process-tensor-construction} are deliberately empty. In the next section we develop the quantum instruments that may occupy those slots; only afterwards will we assemble them into complete multi-time experiments and contract them with the process tensor.